\textbf{Phần 1. Trắc nghiệm lựa chọn}
\setcounterpageref{ex}{0}

\begin{ex} %[3D8N46]
Điền dấu thích hợp: $4\,520 \dots 4\,250$.
\choice
{$<$}
{\True $>$}
{$=$}
{Không so sánh được}
\loigiai{
Hàng nghìn bằng nhau, so sánh hàng trăm: $5 > 2$ nên $4520 > 4250$.\\
\textbf{Đáp án của câu hỏi là B}
}
\end{ex}

\begin{ex} %[3D8N46]
Điền dấu thích hợp: $8\,999 \dots 9\,000$.
\choice
{$>$}
{$=$}
{\True $<$}
{Không so sánh được}
\loigiai{
$8999 < 9000$.\\
\textbf{Đáp án của câu hỏi là C}
}
\end{ex}

\begin{ex} %[3D8H46]
Số nào sau đây lớn nhất?
\choice
{$5432$}
{\True $5480$}
{$5423$}
{$5408$}
\loigiai{
Số lớn nhất là $5480$.\\
\textbf{Đáp án của câu hỏi là B}
}
\end{ex}

\begin{ex} %[3D8H46]
Số nào sau đây bé nhất?
\choice
{$6780$}
{$6708$}
{$6870$}
{\True $6087$}
\loigiai{
Số bé nhất là $6087$.\\
\textbf{Đáp án của câu hỏi là D}
}
\end{ex}

\begin{ex} %[3D8H46]
Sắp xếp các số $3450, 3540, 3054, 3405$ theo thứ tự từ bé đến lớn:
\choice
{\True $3054 < 3405 < 3450 < 3540$}
{$3540 < 3450 < 3405 < 3054$}
{$3054 < 3450 < 3405 < 3540$}
{$3405 < 3054 < 3450 < 3540$}
\loigiai{
Thứ tự tăng dần: $3054 < 3405 < 3450 < 3540$.\\
\textbf{Đáp án của câu hỏi là A}
}
\end{ex}

\begin{ex} %[3D8H46]
Điền chữ số thích hợp vào ô trống để được $5\,67\dots > 5\,678$:
\choice
{$7$}
{$8$}
{\True $9$}
{$6$}
\loigiai{
Để lớn hơn $5678$, chữ số hàng đơn vị phải là 9.\\
\textbf{Đáp án của câu hỏi là C}
}
\end{ex}

\begin{ex} %[3D8V46]
Trong các số $4\,560; 5\,460; 4\,650; 5\,640$, số lớn nhất là
\choice
{$5460$}
{$4650$}
{$4560$}
{\True $5640$}
\loigiai{
Số lớn nhất là $5640$.\\
\textbf{Đáp án của câu hỏi là D}
}
\end{ex}

\begin{ex} %[3D8V46]
Tìm số tự nhiên $x$ tròn nghìn biết: $4\,500 < x < 6\,500$.
\choice
{$4000$}
{\True $5000$ (hoặc $6000$)}
{$7000$}
{$3000$}
\loigiai{
Số tròn nghìn thỏa mãn là $5000$ hoặc $6000$.\\
\textbf{Đáp án của câu hỏi là B}
}
\end{ex}

\begin{ex} %[3D8V46]
Giá tiền một chiếc áo là $85\,000$ đồng (hoặc $8\,500$ đồng tùy phạm vi), một chiếc mũ là $7\,500$ đồng. Vật nào có giá rẻ hơn?
\choice
{\True Chiếc mũ}
{Chiếc áo}
{Bằng nhau}
{Không so sánh được}
\loigiai{
Vì $7500 < 8500$ nên chiếc mũ rẻ hơn.\\
\textbf{Đáp án của câu hỏi là A}
}
\end{ex}

\begin{ex} %[3D8V46]
Tìm chữ số $x$ biết $7x50 < 7150$:
\choice
{$1$}
{$2$}
{\True $0$}
{$3$}
\loigiai{
Để $7x50 < 7150$ thì $x < 1 \Rightarrow x = 0$.\\
\textbf{Đáp án của câu hỏi là C}
}
\end{ex}

\begin{ex} %[3D8V46]
Dãy số nào sau đây được xếp theo thứ tự giảm dần?
\choice
{\True $9870; 9780; 8970; 8790$}
{$4520; 4530; 4540; 4550$}
{$1234; 2341; 3412; 4123$}
{$5000; 5500; 5050; 5005$}
\loigiai{
Dãy giảm dần: $9870 > 9780 > 8970 > 8790$.\\
\textbf{Đáp án của câu hỏi là A}
}
\end{ex}

\begin{ex} %[3D8C46]
Tìm số tròn chục $x$ biết $2\,450 < x < 2\,470$:
\choice
{$2455$}
{$2470$}
{$2450$}
{\True $2460$}
\loigiai{
Số tròn chục nằm giữa là $2460$.\\
\textbf{Đáp án của câu hỏi là D}
}
\end{ex}

\textbf{Phần 2. Trắc nghiệm đúng sai}
\setcounterpageref{ex}{0}

\begin{ex} %[3D8H46]
Xét tính đúng sai khi so sánh các số:
\choiceTF
{\True $3456 < 3465$.}
{\True $7000 > 6999$.}
{$5000 + 400 > 5401$.}
{\True $5000 + 400 < 5401$ (vì $5400 < 5401$).}
\loigiai{
\begin{itemize}
    \item $3456 < 3465$: Đúng.
    \item $7000 > 6999$: Đúng.
    \item $5400 < 5401$ (không lớn hơn). Ý 3 Sai, Ý 4 Đúng.
\end{itemize}
\begin{itemize}
    \item a) \textbf{Đúng}.
    \item b) \textbf{Đúng}.
    \item c) \textbf{Sai}.
    \item d) \textbf{Đúng}.
\end{itemize}
}
\end{ex}

\begin{ex} %[3D8V46]
Cho 4 số: $6\,120; 6\,210; 6\,021; 6\,201$. Xét tính đúng sai:
\choiceTF
{\True Số lớn nhất là $6\,210$.}
{\True Số bé nhất là $6\,021$.}
{Có 3 số lớn hơn $6\,200$.}
{\True Có 2 số lớn hơn $6\,200$ (là $6\,201$ và $6\,210$).}
\loigiai{
\begin{itemize}
    \item Lớn nhất: $6210$. Đúng.
    \item Bé nhất: $6021$. Đúng.
    \item Lớn hơn 6200 chỉ có 6201 và 6210 (2 số). Ý 3 Sai, Ý 4 Đúng.
\end{itemize}
\begin{itemize}
    \item a) \textbf{Đúng}.
    \item b) \textbf{Đúng}.
    \item c) \textbf{Sai}.
    \item d) \textbf{Đúng}.
\end{itemize}
}
\end{ex}

\textbf{Phần 3. Trả lời ngắn}
\setcounterpageref{ex}{0}

\begin{ex} %[3D8H46]
Tìm số lớn nhất trong các số: $4\,560; 4\,650; 4\,506; 4\,605$.
\shortans{4650}
\loigiai{
\textbf{Đáp số:} $4650$\\
$4650$.
}
\end{ex}

\begin{ex} %[3D8H46]
Tìm số bé nhất trong các số: $8\,120; 8\,012; 8\,201; 8\,102$.
\shortans{8012}
\loigiai{
\textbf{Đáp số:} $8012$\\
$8012$.
}
\end{ex}

\begin{ex} %[3D8V46]
Tìm chữ số $x$ để: $3\,4x5 < 3\,415$.
\shortans{0}
\loigiai{
\textbf{Đáp số:} $0$\\
$x = 0$.
}
\end{ex}

\begin{ex} %[3D8V46]
Tìm số tròn trăm $x$ biết: $3\,800 < x < 4\,000$.
\shortans{3900}
\loigiai{
\textbf{Đáp số:} $3900$\\
$x = 3900$.
}
\end{ex}

\begin{ex} %[3D8V46]
Có bao nhiêu số tự nhiên nằm giữa hai số $2\,998$ và $3\,002$?
\shortans{3}
\loigiai{
\textbf{Đáp số:} $3$\\
Các số là $2999, 3000, 3001$ (3 số).
}
\end{ex}

\begin{ex} %[3D8C46]
Tìm số có 4 chữ số lớn nhất có các chữ số khác nhau và tổng các chữ số bằng 10.
\shortans{7210}
\loigiai{
\textbf{Đáp số:} $7210$\\
Để lớn nhất ta chọn hàng nghìn là 7, trăm là 2, chục là 1, đơn vị là 0 ($7 + 2 + 1 + 0 = 10$). Số đó là $7210$.
}
\end{ex}
